\section{What remains, ranked}\label{sec:outlook}

The decomposition closes some questions and opens others. We list what is
closed first, so that it is not redone, and then what is open, in the order we
would take it.

\subsection{Closed}

\begin{enumerate}
\item The comb is not a Loewner phenomenon, in either identification of the
variables, for single- or multi-tip chains, deterministic or L\'evy-driven
(Section~\ref{sec:nocomb}).
\item The rational correction carries no arithmetic and its residues are not
data (Proposition~\ref{prop:forced}); the inherited note's expectation that a
realization would fail ``at two numbers'' is withdrawn.
\item Structural positivity of the Markov part bounds nothing about
$\norm{V_{\omega,L}}$ (Section~\ref{sec:decomposition}); the cancellation
between the positive part and its EMA is the content.
\item The far-right-line Fourier realization of the transfer is superseded for
numerical purposes by the elementary assembly (Section~\ref{sec:assembly}).
\item \textbf{The contraction side is finished.} The assembly is registered and
replayed in double precision (Appendix~\ref{app:numerics}); the first-order law
is verified at $L=\log3,\log5,\log7$ over fourteen orders of magnitude of margin,
with the deviation identified as the Galerkin term of
Proposition~\ref{prop:galerkin} and measured at each horizon
(Table~\ref{tab:horizons}); the nested defect is carried to $N'=192$ and
overshoots as \eqref{eq:nested-limit} requires (Table~\ref{tab:nested}); the
cumulative Cayley coordinate is computed, its congruence checked, and its parity
in $\omega$ proved (Section~\ref{sec:cayleycoord}); and the phase-matched Gram
control is compared with the transfer's own margin in one basis
(Section~\ref{sec:lattice}).
\end{enumerate}

\subsection{Open}

\begin{enumerate}
\item \textbf{The second-order coefficient.} The exact compression has
$\lambda_{\min}(D)/2\omega=m_L-\omega m_L^2+c_2\omega^2+O(\omega^3)$ and
$\lambda_{\min}(P)=m_L+c_2'\omega^2+O(\omega^4)$, with the difference
$c_2'-c_2=0.0362\,m_L$ measured at $L=\log3$, $N=24$ \eqref{eq:cayleygap}. Here
$c_2$ involves $2Q^2+[Q,A]$ by \eqref{eq:secondorder} and $c_2'$ involves
$\re(G^3)$ and the $\omega^3$ term of the generator. Computing either against
the measurement tests the expansion at the next order, and is the only place so
far where $A$ --- the compressed Hilbert transform of the interval together with
the odd part of the comb --- enters a number that has been measured. Within
reach of the present instrument, and first.
\item \textbf{The Lax--Phillips dictionary at $\omega=\frac12$}, object by
object: which subspace the compression $P_L$ is in their energy form, what the
flux identity \eqref{eq:flux} is there, and what their semigroup estimate says
about the contraction margin at the endpoint. Then which of their objects the
$\omega$-deformation moves and which it cannot reach.
\item \textbf{The Beta law from a radial chain.} Whether a conformal-radius or
hitting law of radial Loewner evolution driven by a Bessel process of dimension
$b$ has the additivity fraction of Proposition~\ref{prop:bessel} as a squared
radius ratio. Remark~\ref{rem:betavar} sharpens the target: the archimedean
delay is exactly $\tau=-\frac12\log U$ with $U\sim\mathrm{Beta}(\frac a2,\omega)$,
so what a chain would have to produce is $e^{-2\tau}$ Beta-distributed. If yes,
the archimedean factor has a Loewner realization and its quarter is a Bessel
dimension; if no, the archimedean factor is a Bessel-process statement and not a
Loewner one, and the investigation's name records its origin rather than a
direction.
\item \textbf{The comb in the Bost--Connes algebra.} $C_\omega$ is the image of
$\zeta(s-\omega)/\zeta(s+\omega)$ under $n^{-s}\mapsto\mu_n$. What the KMS
structure of that algebra says about the compression $P_LC_\omega P_L$, whose
norm grows like $e^{bL}$, and whether the subtraction $2\EMA_b$ has a meaning
there. A reading task with a small chance of a real statement.
\item \textbf{The criticality in filter terms.} The transfer is a lossless
filter with impulse response $\kh_\omega-2\EMA_b[\kh_\omega]$, $\kh_\omega\geq0$,
and RH is its stability at every $\omega>0$. A clean statement of why no
passivity theorem for that class can avoid the location of the poles --- what
property of $\kh_\omega$ beyond positivity would be needed --- would be the
honest end of this line.
\end{enumerate}

\subsection{Status of every statement}

\emph{Written proof, unconditional:} Propositions~\ref{prop:unimodular},
\ref{prop:factor}--\ref{prop:blaschke}, \ref{prop:bessel}, \ref{prop:farfield},
\ref{prop:analytic}--\ref{prop:tail}, \ref{prop:congruence}--\ref{prop:cayleygalerkin},
Lemmas~\ref{lem:multiplicative} and~\ref{lem:rotation}, Remark~\ref{rem:betavar},
and the identities of Section~\ref{sec:endpoint}.
\emph{Quoted from \cite{WilsonLines}:} Proposition~\ref{prop:dichotomy}, the
flow identities, and the Cayley coordinate \eqref{eq:cayley} with its
positive-real criterion. \emph{Registered computation, replayed
byte-identically:} the double-precision part of Section~\ref{sec:results} and
Section~\ref{sec:cayleycoord}, 53 cases (Appendix~\ref{app:numerics}).
\emph{Labelled numerical computation (unregistered):} the remainder of
Sections~\ref{sec:results}, \ref{sec:cayleycoord} and~\ref{sec:lattice}, and
Appendix~\ref{app:numerics}. \emph{Reading:} the Lax--Phillips remark of
Section~\ref{sec:endpoint}; the Bost--Connes identification beyond the algebra
written out; the inferred $\norm{(I-P_{24})Af_{24}}^2=1.3265\,m_L^{(24)}$ of
Section~\ref{sec:results}, which combines a measured slope with a computed term;
and the reading of Section~\ref{sec:lattice}. \emph{Superseded:} the two
smallest-$\omega$ margin rows of \cite[Table~3.1]{ContractionNote}, which lay
beneath that note's own quadrature floor, and its $1/N'$ reading of the nested
excess \cite{RegisteredNote}. \emph{Not claimed:} a realization of the transfer,
a positivity statement, contraction at any horizon where $m_L>0$ is not already
known, convergence of the sine basis, or anything about the zeros.
